\section{Accuracy allocation with endogenous successor cost}\label{sec:coupled89}
The separable allocation in Proposition~\ref{prop:allocation86} applies when its coefficients are fixed independently of the allowances being optimized. A dense successor circuit generally has a size that depends on its own allowance. Charging that size at the preceding date couples the two choices. We now optimize this dependency explicitly.

Let $\eta=(\eta_0,\ldots,\eta_{T-1})$ be positive dated residual allowances. Primitive analysis of a specified stored-continuation implementation supplies the resource upper model
\begin{equation}\label{eq:coupled-work89}
 F(\eta)=F_0+\sum_{j=1}^{J}c_j\prod_{t<T}\eta_t^{-a_{jt}},
 \qquad c_j>0,\quad a_{jt}\ge0.
\end{equation}
The coefficients, exponents and setup term $F_0\ge0$ are fixed before choosing $\eta$. A term $c_t\eta_t^{-p_t}\eta_{t+1}^{-s_{t+1}}$ charges date-$t$ label acquisition and evaluation of its stored successor. Compilation, terminal evaluation and storage-related work may supply other terms. An operation ledger is scalarized only at declared nonnegative resource prices; the model does not estimate monetary prices or energy use.

The feasible accuracy set is
\begin{equation}\label{eq:coupled-budget89}
 \mathcal A=\{\eta:0<\eta_t\le\bar\eta_t,\quad
                    \sum_{t<T}w_t\eta_t\le\varepsilon\},
 \qquad w_t>0,\quad\bar\eta_t>0,\quad\varepsilon>0.
\end{equation}
For the original expected-cost sandwich, $w_t=B_t$. The caps specify where the implementation's error and resource bounds hold. A coordinate with no resource-dependent term can be fixed separately; assume here that every coordinate has a positive exponent in at least one term.

\begin{theorem}[Coupled accuracy optimization with an exact certificate]\label{thm:coupled89}
The upper-model minimization $F_* =\min_{\mathcal A}F(\eta)$ has a unique positive minimizer. It is a convex program in $\eta$ and, after $\eta_t=\exp(x_t)$, also has a convex exponential representation. Standard geometric-programming methods can propose an allocation, but their numerical status is not its certificate.

For any feasible $\eta^0$, write $g=\nabla F(\eta^0)$. Let $\lambda\ge0$ and $\nu_t\ge0$ satisfy
\begin{equation}\label{eq:coupled-certificate89}
 g_t+\lambda w_t+\nu_t\ge0\qquad(t<T).
\end{equation}
Then the independently checkable bounds are
\begin{align}\label{eq:coupled-dual89}
 D(\eta^0,\lambda,\nu)
 &=F(\eta^0)-g\cdot\eta^0-\lambda\varepsilon
                  -\sum_t\nu_t\bar\eta_t,\notag\\
 D(\eta^0,\lambda,\nu)&\le F_*\le F(\eta^0).
\end{align}
At the optimizer, the Karush--Kuhn--Tucker multipliers make these bounds equal. For any candidate, the choice $\nu=0$ and
$\lambda=\max_t\{-g_t/w_t\}$ supplies a valid, possibly loose, lower bound. With integer exponents and rational inputs and candidates, every check is rational arithmetic. For other exponents, certified enclosures of the function and derivatives are required.

Put $P=\max_j\sum_ta_{jt}$. Rounding each allowance down by less than a factor two preserves feasibility and gives
\[
 F(\eta^{\rm dyad})-F_0
 \le2^P\{F(\eta^0)-F_0\}.
\]
The exponent is the total degree of a coupled term, not merely the exponent of the date at which that term is paid.
\end{theorem}

The proof in Supplement~\ref{supp:coupled89} includes attainment, strict convexity and the certificate. Convex duality and geometric programming are established techniques \citep{boyd2004,boyd2007gp}; their role here is to preserve the successor dependency in the NBO's original accuracy account. The certificate bounds the optimum of an explicitly stated resource upper model. It is not a lower bound on the work of every algorithm solving the economic problem.

\begin{corollary}[A constructive horizon bound with coupled evaluation]\label{cor:coupled-horizon89}
Suppose $w_t=1$, $\varepsilon/T\le\min\{1,\min_t\bar\eta_t\}$, $J\le KT$, $c_j\le C$, and the total degrees are bounded by a fixed $P$. Then a feasible stored-continuation implementation satisfies
\[
 F_*\le F_0+KC\,T^{P+1}\varepsilon^{-P}.
\]
For a dense successor account with label exponent $p$ and evaluation exponent $s$, the adjacent-date term has degree $p+s$. Thus a valid polynomial horizon account charges $p+s$, rather than silently holding the successor cost constant. Worsening regularity, growing cover dimension, bit lengths and setup remain explicit inputs; none is bounded by this corollary without its stated hypotheses.
\end{corollary}

\paragraph{A priced two-date allocation.}
Consider the exact upper model
\[
 F(\eta_0,\eta_1)=\eta_0^{-1}
             +(\eta_0\eta_1)^{-1}+2\eta_1^{-2},
 \qquad\eta_0+\eta_1\le\tfrac12,\quad\bar\eta_t=1.
\]
The equal allocation costs $52$. The current numerical appendix gives a rational feasible allocation, its independently recomputed lower certificate and its gap. The example isolates the accuracy dependency algebraically; its resource units are not fitted runtime measurements. It also provides a regression test that rejects a negative multiplier, an infeasible budget or a forged lower bound. In the original economy the same construction applies only after its coefficients have been derived from the actual label, action, integration and stored-circuit accounts.

\paragraph{Policy and resource certificates.}
Qualification determines which policy outputs meet the economic accuracy requirement. The coupled program determines how a specified implementation distributes the corresponding dated allowances. These are different optimizations. A tight resource certificate cannot replace an original-law policy certificate, and a tight policy certificate does not determine hardware time. Keeping both accounts explicit permits constructive improvement without a claim of dimension-free or universally faster computation.
